\chapter{Marco Te\'orico}

\section{La diabetes y su impacto en la salud}

\subsection{Qu\'e es la diabetes}
\subsubsection{Descripci\'on de la enfermedad}
\subsubsection{Consecuencias en la salud p\'ublica}
\subsubsection{Importancia de detectarla a tiempo}

\subsection{Tipos principales de diabetes}
\subsubsection{Diabetes juvenil (Tipo 1)}
\subsubsection{Diabetes mellitus (Tipo 2)}
\subsubsection{Diabetes gestacional}

\section{Factores y Datos M\'edicos Relacionados con la Diabetes}

\subsection{Factores que aumentan el riesgo de diabetes}
\subsubsection{Factores gen\'eticos}
\subsubsection{Edad y estilo de vida}
\subsubsection{Antecedentes familiares}

\subsection{Indicadores m\'edicos para detectar diabetes}
\subsubsection{Glucosa en sangre}
\subsubsection{Hemoglobina glucosilada (HbA1c)}
\subsubsection{Presi\'on arterial}
\subsubsection{Colesterol HDL}
\subsubsection{\'Indice de masa corporal}

\subsection{Evaluaci\'on del riesgo hereditario}
\subsubsection{Influencia gen\'etica en la diabetes}
\subsubsection{Funci\'on Pedigree Diabetes Function}
\subsubsection{Interpretaci\'on del riesgo familiar}

\section{Formas Actuales de Detectar la Diabetes}

\subsection{M\'etodos m\'edicos tradicionales}
\subsubsection{Pruebas cl\'inicas}
\subsubsection{Evaluaci\'on mediante historial m\'edico}

\subsection{Problemas de los m\'etodos tradicionales}
\subsubsection{Diagn\'ostico tard\'io}
\subsubsection{Dependencia del an\'alisis manual}
\subsubsection{Dificultad para analizar m\'ultiples factores}

\section{Uso de la Inteligencia Artificial en la Salud}

\subsection{Aplicaci\'on de la Inteligencia Artificial en salud}
\subsubsection{Qu\'e es la Inteligencia Artificial}
\subsubsection{Uso de IA en el \'area m\'edica}
\subsubsection{IA en el diagn\'ostico de enfermedades}

\subsection{Machine Learning en enfermedades cr\'onicas}
\subsubsection{Qu\'e es el aprendizaje autom\'atico}
\subsubsection{Modelos predictivos de salud}
\subsubsection{Ventajas en el an\'alisis de datos cl\'inicos}

\subsection{Redes Neuronales en el diagn\'ostico m\'edico}
\subsubsection{C\'omo funcionan}
\subsubsection{Uso en predicci\'on de enfermedades}
\subsubsection{Aplicaci\'on en detecci\'on de diabetes}


\section{Manejo y An\'alisis de Datos M\'edicos}

\subsection{Registros electr\'onicos de salud}
\subsection{Procesamiento y an\'alisis de datos cl\'inicos}
\subsection{Integraci\'on de variables m\'edicas en modelos predictivos}

\section{Sistemas Inteligentes para Detectar la Diabetes}

\subsection{Sistemas de apoyo para m\'edicos}

\subsection{Sistemas predictivos para el diagn\'ostico temprano}
\subsubsection{Funcionamiento del modelo predictivo}
\subsubsection{Variables utilizadas en la predicci\'on}
\subsubsection{Proceso de entrenamiento del modelo}
\subsubsection{Evaluaci\'on del rendimiento del modelo predictivo}

\subsection{Sistemas predictivos en la prevenci\'on de la diabetes}

\subsection{Importancia de la IA en la detecci\'on m\'edica}

\section{Consideraciones \'eticas en el uso de Inteligencia Artificial en la salud}

\subsection{Protecci\'on de datos personales}
\subsection{Uso responsable de la inteligencia artificial}
\subsection{Transparencia y confiabilidad de los sistemas}